%% annex/p07/calculi-and-cut-elimination-supplement.tex
% Necessity: the exact inherited proof sketches do not establish the distinct
% registered-calculus soundness and presentation-translation claims.
% Definition: the opening two items retain the historical lines 5--35 sketch
% material; the later items are new complete proofs of LK--ND equivalence and
% canonical-LK completeness.
% Standing: it is separately attributed editorial evidence, not source proof
% text, a source-status upgrade, or a completed review determination.
% Nonredundancy: its body labels point only to the distinct registered-calculus
% additions and not to the exact source theorem or its proof sketch.
% Use: annex/p07.tex routes it once as editorial annex material, not restored
% source, with reciprocal body destinations for every item.

\section{Registered-calculus additions}

% BEGIN HISTORICAL SLICE 5-35
\begin{annexitem}{annex:p7.calculi.1}{thm:p07-lkd-soundness}{Soundness of the defined-negation sequent calculus}
\sketchnote{the individual semantic calculations for every structural,
propositional, and quantifier rule; in particular, the substitution-lemma
calculation and the two eigenvariable cases.}
Induct on the height of a
\(\mathsf{LK}_{\mathrm d}\)-derivation.  The initial and structural rules
preserve the semantic reading of the sequent.  Each propositional case unfolds
the truth clause of its principal connective.  The quantifier cases use the
first-order substitution lemma; freshness of the eigenvariable permits the
assignment variation in \(\forall R\) and \(\exists L\).  Cut composes the
two validity implications through its displayed formula.
\end{annexitem}

\begin{annexitem}{annex:p7.calculi.2}{thm:p07-lk-presentation-equivalence}{Canonical/defined-negation LK equivalence}
\sketchnote{contexting uncontexted identity and its converse; admissibility
or derivability of the translated primitive-negation rules; both directions
between shared-context and split-context implication-left using weakening and
contraction; the remaining propositional cases; and all four
quantifier/eigenvariable cases.  No cut-free, height, or size preservation is
established.}
Translate formulas occurrencewise by \(d\).  In the canonical-to-defined
direction, induct on the canonical derivation: simulate contexted initial
sequents by uncontexted identity plus weakenings, expand primitive negation as
implication to \(\bot\), and split shared implication-left contexts using the
registered structural rules.  In the reverse direction, use canonical
contexted identity and structural rules to simulate the uncontexted and
split-context presentation.  Homomorphic connective and quantifier clauses
then preserve each translated endsequent.  The omitted rule-by-rule
verification is why the result is a sketch.
\end{annexitem}

% END HISTORICAL SLICE 5-35

\begin{annexitem}{annex:p7.calculi.3}{thm:p07-lk-nd-equivalence}{Derivability equivalence of canonical LK and classical natural deduction}
\begin{proof}
Write \(D_{\Delta}=\bigvee\Delta\), with the convention and fixed association
of Definition~\ref{def:p07-sequent-nd-translation}; in particular,
\(D_{\varnothing}=\bot\).  Repeated \(\lor I\) and \(\lor E\) give two finite
natural-deduction constructions that will be used below: if the support of
\(\Delta\) is contained in that of \(\Pi\), then
\(D_{\Delta}\vdash D_{\Pi}\), and a derivation of
\(D\lor A\lor B\) can be reassociated to a derivation of
\(D\lor(A\lor B)\).  When the smaller support is empty the first construction
is \(\bot E\).  These are derivations, not identifications of differently
associated formulas.

The sheet requires every discharge to be nonvacuous.  Whenever a construction
below has a derivation of \(C\) under which a displayed assumption \(A\) would
otherwise be unused, insert the finite detour
\[
 A\qquad C\quad\xRightarrow{\land I}\quad A\land C
 \quad\xRightarrow{\land E_2}\quad C.
 \tag{1}
\]
The marked occurrence of \(A\) then lies above the discharge.  Thus the usual
weakening admissibility is available without violating the exact nonvacuous
rule sheet.

First translate an LK derivation of \(\Gamma\Rightarrow\Delta\) into a
natural-deduction derivation of \(D_{\Delta}\) from
\(\operatorname{supp}(\Gamma)\), by induction on the LK derivation.  Identity
uses its displayed assumption followed by the finite-disjunction inclusion;
\(\bot L\) uses \(\bot E\).  Left weakening uses (1), right weakening uses
finite-disjunction inclusion, and either contraction changes only a
multiplicity forgotten by support.  For Cut, the translated left premise gives
\(D_{\Delta}\lor A\), and the translated right premise gives
\(D_{\Sigma}\) under assumption \(A\).  Eliminate the first disjunction:
its left branch injects \(D_{\Delta}\) into
\(D_{\Delta,\Sigma}\), while its right branch applies the second translation
and injects \(D_{\Sigma}\).  This derives the translated cut conclusion.

The left logical rules translate by their matching eliminations.  For
\(\land L\), obtain both conjuncts and use the translated premise; for
\(\lor L\), apply \(\lor E\) to the two translated premises; for
\(\to L\), eliminate the translated formula \(D\lor A\): the \(D\)-branch
is already a required succedent, while the \(A\)-branch applies the open
assumption \(A\to B\), then the translated \(B\)-premise.  The rule
\(\neg L\) is the same case with contradiction and \(\bot E\).  The rules
\(\forall L\) and \(\exists L\) use respectively \(\forall E\) and
\(\exists E\); the LK eigenvariable condition for \(\exists L\) is exactly the
condition needed by the natural-deduction elimination.

For the right rules, \(\lor R\) is reassociation and injection, and
\(\exists R\) applies \(\exists I\) in the witness branch.  The translation of
\(\land R\) has premises \(D\lor A\) and \(D\lor B\).  Eliminate the first;
its \(D\)-branch is done, and in its \(A\)-branch eliminate the second, using
its \(D\)-branch directly and combining \(A,B\) by \(\land I\) in the other.
This yields \(D\lor(A\land B)\).

Three right-rule cases use classical RAA.  For \(\neg R\), the induction
hypothesis gives \(D\) under assumption \(A\).  Assume
\(\neg(D\lor\neg A)\).  From it derive \(\neg D\): an assumed \(D\) injects
into the forbidden disjunction.  Also derive \(\neg\neg A\): an assumed
\(\neg A\) injects into the same forbidden disjunction.  RAA gives \(A\), the
translated premise gives \(D\), and \(D\) contradicts \(\neg D\).  A final
RAA yields \(D\lor\neg A\).

For \(\to R\), the induction hypothesis gives \(D\lor B\) under assumption
\(A\).  Assume \(\neg(D\lor(A\to B))\) and derive \(\neg D\).  Under
assumption \(A\), eliminate \(D\lor B\); its \(D\)-branch gives
contradiction and then \(B\), and its other branch already gives \(B\).
Use (1) if needed and discharge \(A\), obtaining \(A\to B\), which contradicts
the outer assumption after injection.  RAA gives the required disjunction.

For \(\forall R\), the translated premise is \(D\lor A(y)\), where \(y\) is
free in neither the conclusion sequent nor its contexts.  Assume
\(\neg(D\lor\forall xA(x))\) and derive \(\neg D\).  Eliminating
\(D\lor A(y)\) gives \(A(y)\), since the \(D\)-branch yields contradiction
and then \(A(y)\).  The variable \(y\) is free in no remaining open assumption,
including the temporary negation, so \(\forall I\) gives \(\forall xA(x)\).
Injection contradicts the temporary assumption, and RAA completes the case.
The \(\neg R\), \(\to R\), and outer RAA assumptions are each used in the
displayed contradiction; hence their discharges are nonvacuous.  This exhausts
all rules of equality-free \texttt{sheet:lk}, including the four quantifier
rules and empty \(D\).

Conversely, induct on an equality-free derivation on
\hyperref[sheet:nd-fol]{\texttt{sheet:nd-fol}} to obtain a canonical-LK
derivation of \(\Gamma\Rightarrow A\) from every judgment
\(\Gamma\vdash A\).  The assumption case is contexted identity.  Contexts in
two- and three-premise natural-deduction rules are aligned by LK weakening and
their duplicate formula occurrences are removed by contraction.
The rules \(\land I\), \(\to I\), \(\neg I\), \(\forall I\), and
\(\exists I\) translate directly by the corresponding right rule (for
\(\neg I\), first cut the premise succedent \(\bot\) against
\(\bot\Rightarrow\), then use \(\neg R\)).  The eliminations
\(\land E_i\), \(\to E\), and \(\neg E\) cut the translated premises against,
respectively, the derived sequents
\[
 A\land B\Rightarrow A,\qquad
 A,A\to B\Rightarrow B,
 \qquad
 A,\neg A\Rightarrow\bot.
\]
The first is \(\land L\) over identity; the second is the \(\to L\) kernel
shown in Figure~\ref{fig:p07-sequent-nd-correspondence}; the third is
\(\neg L\) over contexted identity.  The rule \(\bot E\) cuts with the
initial sequent \(\bot\Rightarrow A\).

For \(\lor I_i\), weaken the succedent by the other disjunct and apply
\(\lor R\).  For \(\lor E\), align the two minor-premise contexts, apply
\(\lor L\) to obtain \(A\lor B,\Delta,\Sigma\Rightarrow C\), and cut the
translated major premise.  Universal elimination cuts against
\(\forall xA\Rightarrow A[t/x]\), obtained by \(\forall L\) over identity;
universal introduction is \(\forall R\), with the same freshness condition.
Existential introduction is \(\exists R\); existential elimination first
applies \(\exists L\) to the translated minor premise and then cuts the
translated major premise.  The eigenvariable condition is unchanged.
Alpha-conversion is admissible by induction on an LK derivation, consistently
renaming bound variables in every formula instance.  The two equality rules
are absent by the equality-free hypothesis.

For RAA, the translated premise is
\(\Gamma,\neg A\Rightarrow\bot\).  Cut it with \(\bot\Rightarrow\) to obtain
\(\Gamma,\neg A\Rightarrow\); apply \(\neg R\) to get
\(\Gamma\Rightarrow\neg\neg A\).  A cut-free derivation of
\(\neg\neg A\Rightarrow A\) is obtained by applying \(\neg R\) to identity
\(A\Rightarrow A\), producing \(\Rightarrow A,\neg A\), and then applying
\(\neg L\) to the displayed \(\neg A\).  Cut with it to obtain
\(\Gamma\Rightarrow A\).

The induction gives
\(\operatorname{supp}(\Gamma)\Rightarrow D_{\Delta}\) from a derivation of
\(D_{\Delta}\) in natural deduction.  Finally, canonical LK derives
\(D_{\Delta}\Rightarrow\Delta\) by induction on the finite support of
\(\Delta\): the empty case is \(\bot L\), the singleton case is contexted
identity, and the inductive case is \(\lor L\) applied to the two smaller
derivations; right weakening then restores any duplicate occurrences in
\(\Delta\).  Cut gives
\(\operatorname{supp}(\Gamma)\Rightarrow\Delta\), and left weakening restores
multiplicities in \(\Gamma\).  This proves the reverse implication and hence
the theorem.
\end{proof}
\end{annexitem}

\begin{annexitem}{annex:p7.calculi.4}{thm:p07-canonical-lk-completeness}{Finite-sequent completeness of canonical LK}
\begin{proof}
Soundness is exactly Theorems~\ref{thm:prop-lk-soundness} and
\ref{thm:fol-lk-soundness}.  For completeness, suppose first that the finite
propositional sequent \(\Gamma\Rightarrow\Delta\) is valid.  Then
\(\operatorname{supp}(\Gamma)\entails D_{\Delta}\), including the case
\(D_{\Delta}=\bot\).  The proved natural-deduction completeness theorem
\ref{thm:prop-completeness} gives
\(\operatorname{supp}(\Gamma)\vdash D_{\Delta}\), and the reverse translation
in Annex~\ref{annex:p7.calculi.3} gives the required LK derivation.

The first-order proof is the same.  The semantic condition on the sequent is
precisely
\(\operatorname{supp}(\Gamma)\entails D_{\Delta}\).  Under Choice,
Theorem~\ref{thm:fol-completeness} supplies a natural-deduction derivation; for
a finite or countable signature use
Corollary~\ref{cor:fol-completeness-countable}.  The equality-free
specialization of those complete Henkin proofs omits \(=I\), \(=E\), the
equality-congruence lemma, and the quotient of terms: the term structure uses
literal terms, while the connective, quantifier, consistency, witness, and
truth-lemma inductions are unchanged.  It therefore returns a derivation in
the equality-free restriction of \texttt{sheet:nd-fol}.  Annex~\ref{annex:p7.calculi.3}
translates that derivation to canonical LK.  The construction uses Cut in the
translation, so no cut-free strengthening follows.
\end{proof}
\end{annexitem}
